\begin{frame}{왜 $\gamma < 1$ 이면 수렴하는가 (증명 개요)}
  벨만 연산자를
  \[
  \mathcal{B}^\pi V(s) = R(s, \pi(s)) + \gamma\, \mathbb{E}_{s'}[V(s')]
  \]
  라고 쓰자. 두 값 함수 $V, W$ 에 대해 $\gamma < 1$ 이면:
  \[
  \big\|\mathcal{B}^\pi V - \mathcal{B}^\pi W\big\|_\infty
  \;\leq\; \gamma\, \|V - W\|_\infty
  \]
  ($\|\cdot\|_\infty$ = 최대 절댓값, ``어떤 상태에서 가장 많이
  다르냐'')
  \vskip0.4em
  \begin{block}{수축 연산 + 바나흐 고정점 정리}
    즉 \kb{벨만 연산자는 $\gamma$ 를 계수로 \textbf{줄이는 수축
    (squeeze) 연산}. \kb{바나흐 고정점 정리(Banach fixed-point
    theorem)}에 의해, 수축 연산은 \textbf{고정점이 딱 하나} 존재하고,
    \textbf{어떤 초기값에서 출발해도} 그 고정점으로 수렴한다.
    }\vskip0.2em
    {\small 이 한 줄의 정리가 ``가치 반복이 항상 (고유한) 답으로
    수렴한다''는 \textbf{보장의 전부}.}
  \end{block}
  \vskip0.4em
  \begin{itemize}
    \item \kb{$\gamma = 1$ 이면} 수축 상수가 1이 되어 이 보장이
      \textbf{사라진다} --- 수렴하지 않거나, 유한 에피소드여도
      경계가 무한대가 되면 발산할 수 있다.
    \item 3.2절에서 ``$\gamma < 1$ 이면 수렴''이라 말한 것이,
      여기서 \textbf{수학적으로 왜 그런가}가 보인다.
  \end{itemize}
  \vskip0.3em
  \begin{block}{고정점의 유일성}
    $\gamma < 1$ 이면 벨만방정식을 만족하는 $V^\pi$ 는
    \textbf{정확히 하나뿐} --- ``이 정책의 가치함수''라는 개념이
    \textbf{의미가 있는(well-defined)} 것의 근거.
  \end{block}
\end{frame}
